\chapter{Polynomial Transposes of Relative Mealy Actions and Guarded Program Structure}
\label{ch:polynomial-relative-mealy-actions}

\section{Overview}
\label{sec:poly-relative-overview}

The previous chapter developed relative Mealy dynamic proarrow logic from the
Eilenberg--Moore semantics of internal categories regarded as monads in
\(\mathbf{Span}(\mathcal E)\).  A Mealy morphism
\[
\Phi=(P,\delta,\omega):\mathbb A\rightsquigarrow\mathbb B
\]
and a right \(\mathbb B\)-action
\[
Y=(Y\xrightarrow{r}B_0,\sigma)
\]
determine a restricted right \(\mathbb A\)-action
\[
\Phi^\ast Y
\]
on the relative state object
\[
S_{\Phi,Y}:=P\times_{q,B_0,r}Y.
\]
Its action map is
\[
\alpha_{\Phi,Y}:
A_1\times_{t_A,A_0,p\pi_P}S_{\Phi,Y}
\longrightarrow
S_{\Phi,Y},
\]
given in generalized-element notation by
\[
\alpha_{\Phi,Y}(a,x,y)
=
\left(
\delta(a,x),
\sigma(\omega(a,x),y)
\right).
\]
The associated action category is the relative program category
\[
\mathsf{Prog}_{\Phi,Y}
=
\int_{\mathbb A}\Phi^\ast Y.
\]

This chapter explains the polynomial structure that is canonically hidden in
this Eilenberg--Moore action.  The point is not that polynomial functors replace
the span semantics.  Rather, whenever the relevant dependent products exist,
the uncurried action map
\[
A_1\times_{A_0}S_{\Phi,Y}\to S_{\Phi,Y}
\]
may be transposed to a total response profile:
\[
(x,y)
\longmapsto
\left(
a
\longmapsto
\left(
\omega(a,x),
\left(\delta(a,x),\sigma(\omega(a,x),y)\right)
\right)
\right).
\]
Thus the polynomial coalgebra is the dependent-product transpose of the
relative Eilenberg--Moore action, together with its downstream output label.

The guiding slogan is:
\[
\boxed{
\text{span semantics is the uncurried EM execution semantics, while
polynomial semantics is its curried total-response form.}
}
\]

The terminal trace polynomial from the earlier, non-relative presentation is
recovered by taking
\[
Y=\mathbf 1_{\mathbb B}.
\]
In that case
\[
S_{\Phi,\mathbf 1_{\mathbb B}}\cong P,
\]
and the relative response profile reduces to
\[
x
\longmapsto
\left(
a
\longmapsto
(\omega(a,x),\delta(a,x))
\right).
\]

The second theme of the chapter is guarded program structure.  The span
semantics of the previous chapter already supplies tests, guarded primitive
programs, finite choice, sequential composition, path iteration, diamonds, and
boxes.  The polynomial presentation explains one-step guarded modalities as
predicate liftings on total response profiles, pulled back along the relative
Mealy coalgebra.

The conceptual pattern is therefore:
\[
\boxed{
(\Phi^\ast Y,\Lambda_{\Phi,Y})
\quad
\overset{\Pi\text{-transpose}}{\longrightarrow}
c_{\Phi,Y}
\quad
\overset{\text{flatten}}{\longrightarrow}
\mathsf{Prog}_{\Phi,Y}.
}
\]

Here
\[
\Lambda_{\Phi,Y}:
\mathsf{Prog}_{\Phi,Y}\to\int_{\mathbb B}Y
\]
is the downstream comparison functor from the previous chapter.

\section{Standing hypotheses and levels of structure}
\label{sec:poly-relative-hypotheses}

The constructions of this chapter occur at different categorical strengths.

\begin{center}
\[
\begin{array}{c|c|c}
\text{level} & \text{categorical structure} & \text{construction} \\
\hline
0 & \text{pullbacks} &
\Phi^\ast Y,\ \mathsf{Prog}_{\Phi,Y} \\
1 & \text{relevant dependent products} &
\text{polynomial transposes of EM actions} \\
2 & \text{regular images} &
\diamond\text{-modalities and existential response liftings} \\
3 & \text{coherent finite joins} &
\text{finite guarded choice} \\
4 & \text{right adjoints to pullback} &
\Box\text{-modalities and profile boxes} \\
5 & \text{stable countable coproducts} &
\text{path-star iteration} \\
6 & \text{complete fibres and continuity} &
\mu,\nu\text{ fixed-point comparison}
\end{array}
\]
\end{center}

The relative Eilenberg--Moore and program-category constructions only require
pullbacks.  The polynomial layer is additional: it is available when the
specific dependent products appearing below exist.  A sufficient global
hypothesis is that \(\mathcal E\) is locally cartesian closed with finite
products.

For the modal and fixed-point parts of the chapter, we use the regular,
coherent, first-order Heyting, and \(\omega\)-iterative hypotheses introduced
in Chapter~\ref{ch:relative-mealy-program-categories}.

All equalities of spans and predicate transformers are to be read after the
canonical associativity, unit, and Beck--Chevalley identifications.
Equivalently, one may work in the locally posetal quotient of spans by
isomorphism of apexes and in subobject posets.

\section{Polynomial spans as dependent response transformers}
\label{sec:poly-relative-polynomial-spans}

We recall the polynomial notation used throughout the chapter.

For a morphism
\[
f:X\to Y
\]
write
\[
f^\ast:\mathcal E/Y\to\mathcal E/X
\]
for pullback along \(f\),
\[
\Sigma_f:\mathcal E/X\to\mathcal E/Y
\]
for postcomposition with \(f\), and
\[
\Pi_f:\mathcal E/X\to\mathcal E/Y
\]
for the right adjoint to \(f^\ast\), when it exists:
\[
\Sigma_f\dashv f^\ast\dashv \Pi_f.
\]

\begin{definition}[Polynomial span]
\label{def:poly-relative-polynomial-span}
A \textbf{polynomial span} from \(I\) to \(J\) is a diagram
\[
\begin{tikzcd}[column sep=large]
I
&
E
  \arrow[l, "s"']
  \arrow[r, "p"]
&
B
  \arrow[r, "t"]
&
J .
\end{tikzcd}
\]
When the required dependent product exists, it induces a functor
\[
\Sigma_t\Pi_p s^\ast:
\mathcal E/I\to\mathcal E/J.
\]
\end{definition}

Thus a polynomial span is read as the composite
\[
\begin{tikzcd}[column sep=large]
\mathcal E/I
  \arrow[r, "s^\ast"]
&
\mathcal E/E
  \arrow[r, "\Pi_p"]
&
\mathcal E/B
  \arrow[r, "\Sigma_t"]
&
\mathcal E/J .
\end{tikzcd}
\]

In generalized-element notation, if
\[
X\to I
\]
is an object over \(I\), then
\[
\left(\Sigma_t\Pi_p s^\ast X\right)_j
\cong
\sum_{b\in B_j}
\prod_{e\in E_b}
X_{s(e)}.
\]

\begin{remark}[Linear spans]
\label{rem:poly-relative-linear-spans}
An ordinary span
\[
I\xleftarrow{s}R\xrightarrow{t}J
\]
is the special case of a polynomial span in which the shape object is the
position object:
\[
\begin{tikzcd}[column sep=large]
I
&
R
  \arrow[l, "s"']
  \arrow[r, "1_R"]
&
R
  \arrow[r, "t"]
&
J .
\end{tikzcd}
\]
The induced functor is
\[
\Sigma_t\Pi_{1_R}s^\ast
\cong
\Sigma_t s^\ast.
\]
Thus ordinary spans are the linear polynomial case: each execution witness has
exactly one continuation position.  The polynomials below arise by replacing
this single continuation position with a dependent family of input-query
positions.
\end{remark}

\section{Currying ordinary Eilenberg--Moore actions}
\label{sec:poly-relative-currying-em-actions}

Before considering Mealy morphisms, we isolate the basic mechanism for a
right action of an internal category.

Let
\[
\mathbb A=(A_0,A_1,s_A,t_A,e_A,m_A)
\]
be an internal category in \(\mathcal E\).  A right \(\mathbb A\)-action
\[
(X\xrightarrow{p}A_0,\alpha)
\]
has action map
\[
\alpha:
A_1\times_{t_A,A_0,p}X\to X
\]
satisfying
\[
p\alpha=s_A\pi_1.
\]
Equivalently, writing
\[
t_A^\ast X=A_1\times_{t_A,A_0,p}X,
\]
the action map is a morphism over \(A_1\)
\[
t_A^\ast X\to s_A^\ast X.
\]

Assume that the dependent product
\[
\Pi_{t_A}
\]
exists.  Define the \textbf{action-response polynomial} of \(\mathbb A\) by
\[
\mathcal R_{\mathbb A}(X):=\Pi_{t_A}s_A^\ast X.
\]
Fibrewise,
\[
\mathcal R_{\mathbb A}(X)_v
\cong
\prod_{a:u\to v}X_u.
\]

The action map transposes along
\[
t_A^\ast\dashv\Pi_{t_A}
\]
to a morphism over \(A_0\)
\[
\widehat\alpha:
X\to\mathcal R_{\mathbb A}(X).
\]
In generalized-element notation,
\[
\widehat\alpha(x)(a)=\alpha(a,x).
\]

\begin{definition}[Multiplicative action-response coalgebra]
\label{def:poly-relative-multiplicative-action-coalgebra}
A coalgebra
\[
c:X\to\mathcal R_{\mathbb A}(X)
\]
is called \textbf{multiplicative} if, writing
\[
c(x)(a)=x_a,
\]
it satisfies:
\[
c(x)(e_A(px))=x
\]
and, for
\[
a:u\to v,
\qquad
b:v\to w,
\qquad
p(x)=w,
\]
if
\[
c(x)(b)=x_b
\]
and
\[
c(x_b)(a)=x_{ab},
\]
then
\[
c(x)(m_A(a,b))=x_{ab}.
\]
\end{definition}

\begin{proposition}[Actions as multiplicative response coalgebras]
\label{prop:poly-relative-actions-as-response-coalgebras}
Assume the dependent product \(\Pi_{t_A}\) exists.  To give a right
\(\mathbb A\)-action
\[
(X\xrightarrow{p}A_0,\alpha)
\]
is equivalently to give a multiplicative coalgebra
\[
c:X\to\mathcal R_{\mathbb A}(X).
\]
\end{proposition}

\begin{proof}
The adjunction
\[
t_A^\ast\dashv\Pi_{t_A}
\]
identifies maps
\[
t_A^\ast X\to s_A^\ast X
\]
over \(A_1\) with maps
\[
X\to\Pi_{t_A}s_A^\ast X
\]
over \(A_0\).  Under this identification,
\[
\alpha(a,x)
\]
corresponds to
\[
c(x)(a).
\]
The unit and multiplication axioms for the action are exactly the two displayed
multiplicativity equations for \(c\).
\end{proof}

\begin{remark}
\label{rem:poly-relative-first-inevitability}
This proposition is the basic reason polynomial functors appear in this
chapter.  They are not imposed externally.  They are the dependent-product
transposes of Eilenberg--Moore action maps.
\end{remark}

\section{Labelled relative Eilenberg--Moore actions}
\label{sec:poly-relative-labelled-relative-em-actions}

Now let
\[
\Phi=(P,\delta,\omega):\mathbb A\rightsquigarrow\mathbb B
\]
be a Mealy morphism, with carrier span
\[
A_0\xleftarrow{p}P\xrightarrow{q}B_0.
\]
Let
\[
Y=(Y\xrightarrow{r}B_0,\sigma)
\]
be a right \(\mathbb B\)-action.

The restricted right \(\mathbb A\)-action
\[
\Phi^\ast Y
\]
has state object
\[
S_{\Phi,Y}:=P\times_{q,B_0,r}Y
\]
and action
\[
\alpha_{\Phi,Y}(a,x,y)
=
\left(
\delta(a,x),
\sigma(\omega(a,x),y)
\right).
\]

The previous chapter also constructed the downstream comparison functor
\[
\Lambda_{\Phi,Y}:
\mathsf{Prog}_{\Phi,Y}\to\int_{\mathbb B}Y.
\]
On objects,
\[
\Lambda_{\Phi,Y}(x,y)=y,
\]
and on arrows,
\[
\Lambda_{\Phi,Y}(a,x,y)
=
(\omega(a,x),y).
\]

Thus the relative EM action carries more than the next coupled state.  It also
carries the downstream transition produced by the emitted \(\mathbb B\)-arrow.

\begin{definition}[Labelled relative EM action]
\label{def:poly-relative-labelled-relative-em-action}
A \textbf{labelled relative EM action} over a right \(\mathbb B\)-action \(Y\)
consists of:
\begin{enumerate}
\item a right \(\mathbb A\)-action
\[
(X\xrightarrow{p_X}A_0,\alpha);
\]
\item a map
\[
y_X:X\to Y;
\]
\item a map
\[
\lambda:
A_1\times_{t_A,A_0,p_X}X\to B_1
\]
such that, for every admissible pair \((a,z)\),
\[
t_B\lambda(a,z)=r(y_Xz),
\]
and
\[
y_X(\alpha(a,z))=\sigma(\lambda(a,z),y_Xz);
\]
\item the unit and multiplication laws saying that the induced map
\[
(a,z)\longmapsto(\lambda(a,z),y_Xz)
\]
is functorial from the action category \(\int_{\mathbb A}X\) to
\(\int_{\mathbb B}Y\).
\end{enumerate}
\end{definition}

For the labelled action associated to a Mealy morphism,
\[
X=S_{\Phi,Y},
\]
\[
y_X=\pi_Y:S_{\Phi,Y}\to Y,
\]
and
\[
\lambda(a,x,y)=\omega(a,x).
\]

\begin{remark}
\label{rem:poly-relative-need-labelled-action}
The bare restricted action \(\Phi^\ast Y\) remembers only the next coupled
state
\[
(\delta(a,x),\sigma(\omega(a,x),y)).
\]
It need not remember the raw output arrow \(\omega(a,x)\), since different
\(\mathbb B\)-arrows may induce the same transition in \(Y\).  Therefore
output-sensitive guarded modalities are obtained from the labelled relative
EM action, not from the underlying action alone.
\end{remark}

\section{The relative interface object}
\label{sec:poly-relative-interface-object}

Fix a right \(\mathbb B\)-action
\[
Y=(Y\xrightarrow{r}B_0,\sigma).
\]
The natural interface object for the relative polynomial is
\[
J_Y:=A_0\times Y.
\]
Write
\[
\pi_A:J_Y\to A_0,
\qquad
\pi_Y:J_Y\to Y
\]
for the projections.

The relative state object
\[
S_{\Phi,Y}=P\times_{B_0}Y
\]
maps to \(J_Y\) by
\[
j_{\Phi,Y}:S_{\Phi,Y}\to J_Y,
\qquad
j_{\Phi,Y}(x,y)=(p(x),y).
\]

Thus a relative state lies over the interface pair
\[
(p(x),y)\in A_0\times Y.
\]

\begin{remark}[Terminal specialization]
\label{rem:poly-relative-interface-terminal-specialization}
When
\[
Y=\mathbf 1_{\mathbb B}=B_0,
\]
the interface object is
\[
J_Y=A_0\times B_0,
\]
which is the interface object used in the terminal trace presentation.
\end{remark}

\section{Queries and responses}
\label{sec:poly-relative-queries-responses}

We now construct the relative Mealy response polynomial associated to
\[
(\mathbb A,\mathbb B,Y).
\]

\subsection{Input queries}

Define the object of input queries by
\[
Q_Y:=A_1\times_{t_A,A_0,\pi_A}J_Y.
\]
Thus
\[
\begin{tikzcd}[column sep=large,row sep=large]
Q_Y
  \arrow[r, "\overline a"]
  \arrow[d, "\kappa_Y"']
&
A_1
  \arrow[d, "t_A"]
\\
J_Y
  \arrow[r, "\pi_A"']
&
A_0
\end{tikzcd}
\]
is a pullback square.

A generalized element of \(Q_Y\) is written
\[
(a,y),
\]
where
\[
a:u\to v
\]
is an arrow of \(\mathbb A\), and \(y\in Y\).  Its current interface is
\[
(v,y).
\]
Thus
\[
\kappa_Y(a,y)=(t_Aa,y).
\]

There is also a map recording the next input interface:
\[
\lambda_A:Q_Y\to A_0,
\qquad
\lambda_A(a,y)=s_A(a).
\]

\subsection{Output responses}

A response to a query
\[
(a:u\to v,y)
\]
is an arrow of \(\mathbb B\)
\[
\beta:b'\to r(y).
\]
Such an arrow may be consumed by the right \(\mathbb B\)-action \(Y\), giving
the downstream continuation state
\[
\sigma(\beta,y)\in Y.
\]

Define the response object
\[
O_Y
:=
Q_Y\times_{r\pi_Y\kappa_Y,B_0,t_B}B_1.
\]
It fits into a pullback square
\[
\begin{tikzcd}[column sep=large,row sep=large]
O_Y
  \arrow[r, "\overline\beta"]
  \arrow[d, "\rho_Y"']
&
B_1
  \arrow[d, "t_B"]
\\
Q_Y
  \arrow[r, "r\pi_Y\kappa_Y"']
&
B_0 .
\end{tikzcd}
\]
A generalized element is written
\[
(a,y,\beta),
\]
where
\[
t_B(\beta)=r(y).
\]

Define the continuation-interface map
\[
n_Y:O_Y\to J_Y
\]
by
\[
n_Y(a,y,\beta)
=
\left(
s_A(a),
\sigma(\beta,y)
\right).
\]

Thus \(O_Y\) records possible output responses to input queries, and
\(n_Y\) records the interface after the response has been consumed by the
downstream action.

\section{Total response policies and the relative Mealy polynomial}
\label{sec:poly-relative-total-response-policies}

A total output policy at an interface
\[
(v,y)\in J_Y
\]
assigns, to every input arrow
\[
a:u\to v,
\]
an output arrow
\[
\beta_a:b'_a\to r(y).
\]

Formally, define the object of total response policies by the dependent product
\[
K_{\mathbb A,\mathbb B,Y}
:=
\Pi_{\kappa_Y}(O_Y\xrightarrow{\rho_Y}Q_Y).
\]
It is an object over \(J_Y\).  Write
\[
t_{\mathbb A,\mathbb B,Y}:K_{\mathbb A,\mathbb B,Y}\to J_Y
\]
for its structure map.

A generalized element of \(K_{\mathbb A,\mathbb B,Y}\) over
\[
(v,y)
\]
is a policy
\[
\sigma_0:
\{a:u\to v\}
\longrightarrow
\{\beta:b'\to r(y)\}.
\]

Now form the pullback
\[
E_{\mathbb A,\mathbb B,Y}
:=
Q_Y\times_{J_Y}K_{\mathbb A,\mathbb B,Y},
\]
where the maps to \(J_Y\) are
\[
\kappa_Y:Q_Y\to J_Y
\]
and
\[
t_{\mathbb A,\mathbb B,Y}:K_{\mathbb A,\mathbb B,Y}\to J_Y.
\]
Thus
\[
\begin{tikzcd}[column sep=large,row sep=large]
E_{\mathbb A,\mathbb B,Y}
  \arrow[r, "\epsilon_Q"]
  \arrow[d, "p_{\mathbb A,\mathbb B,Y}"']
&
Q_Y
  \arrow[d, "\kappa_Y"]
\\
K_{\mathbb A,\mathbb B,Y}
  \arrow[r, "t_{\mathbb A,\mathbb B,Y}"']
&
J_Y
\end{tikzcd}
\]
is a pullback square.

The counit of
\[
\kappa_Y^\ast\dashv\Pi_{\kappa_Y}
\]
gives an evaluation map
\[
\operatorname{ev}:
E_{\mathbb A,\mathbb B,Y}\to O_Y
\]
over \(Q_Y\).  In generalized-element notation,
\[
\operatorname{ev}(a,y,\sigma_0)=(a,y,\sigma_0(a)).
\]

Define
\[
s_{\mathbb A,\mathbb B,Y}
:=
n_Y\operatorname{ev}
:
E_{\mathbb A,\mathbb B,Y}\to J_Y.
\]
This map records the continuation interface of a queried position in a total
response policy.

\begin{definition}[Relative Mealy response polynomial]
\label{def:poly-relative-mealy-response-polynomial}
The \textbf{relative Mealy response polynomial} associated to
\[
(\mathbb A,\mathbb B,Y)
\]
is the polynomial span
\[
\begin{tikzcd}[column sep=large]
J_Y
&
E_{\mathbb A,\mathbb B,Y}
  \arrow[l, "s_{\mathbb A,\mathbb B,Y}"']
  \arrow[r, "p_{\mathbb A,\mathbb B,Y}"]
&
K_{\mathbb A,\mathbb B,Y}
  \arrow[r, "t_{\mathbb A,\mathbb B,Y}"]
&
J_Y .
\end{tikzcd}
\]
It induces an endofunctor
\[
\mathcal M_{\mathbb A,\mathbb B,Y}:
\mathcal E/J_Y\to\mathcal E/J_Y
\]
given by
\[
\mathcal M_{\mathbb A,\mathbb B,Y}(X)
=
\Sigma_{t_{\mathbb A,\mathbb B,Y}}
\Pi_{p_{\mathbb A,\mathbb B,Y}}
s_{\mathbb A,\mathbb B,Y}^{\ast}X.
\]
\end{definition}

\begin{proposition}[Fibre calculation]
\label{prop:poly-relative-fibre-calculation}
Let
\[
X\to J_Y=A_0\times Y
\]
be an object over the relative interface.  For every generalized interface
\[
(v,y)\in J_Y,
\]
there is a canonical fibrewise description
\[
\mathcal M_{\mathbb A,\mathbb B,Y}(X)_{(v,y)}
\cong
\prod_{a:u\to v}
\sum_{\beta:b'\to r(y)}
X_{(u,\sigma(\beta,y))}.
\]
\end{proposition}

\begin{proof}
A point of
\[
K_{\mathbb A,\mathbb B,Y}
\]
over \((v,y)\) is a total policy assigning to every input arrow
\[
a:u\to v
\]
an output arrow
\[
\beta_a:b'_a\to r(y).
\]
The object
\[
E_{\mathbb A,\mathbb B,Y}
=
Q_Y\times_{J_Y}K_{\mathbb A,\mathbb B,Y}
\]
has, over such a policy, one position for each admissible input arrow
\(a:u\to v\).  The continuation-interface map sends this position to
\[
(u,\sigma(\beta_a,y)).
\]
Therefore
\[
\Pi_{p_{\mathbb A,\mathbb B,Y}}
s_{\mathbb A,\mathbb B,Y}^{\ast}X
\]
assigns to a policy \(\sigma_0\) the family of choices
\[
x_a\in X_{(u,\sigma(\sigma_0(a),y))}
\]
for each input arrow \(a:u\to v\).  Applying
\(\Sigma_{t_{\mathbb A,\mathbb B,Y}}\) regards the result as lying over the
current interface \((v,y)\).

Thus the fibre may first be written as
\[
\sum_{\sigma_0}
\prod_{a:u\to v}
X_{(s_Aa,\sigma(\sigma_0(a),y))}.
\]
We now use the standard dependent distributivity isomorphism available under
the assumed dependent products:
\[
\sum_{\sigma\in\prod_a O_a}
\prod_a X_{a,\sigma(a)}
\cong
\prod_a
\sum_{\beta\in O_a}X_{a,\beta}.
\]
This gives the canonical isomorphism
\[
\sum_{\sigma_0}
\prod_{a:u\to v}
X_{(s_Aa,\sigma(\sigma_0(a),y))}
\cong
\prod_{a:u\to v}
\sum_{\beta:b'\to r(y)}
X_{(u,\sigma(\beta,y))}.
\]
\end{proof}

\begin{remark}
\label{rem:poly-relative-response-polynomial-meaning}
The displayed fibre is the curried one-step response type.  At interface
\((v,y)\), the system must respond to every admissible input arrow
\(a:u\to v\).  For each such \(a\), it chooses an output arrow
\(\beta:b'\to r(y)\) and a continuation state over the interface
\[
(u,\sigma(\beta,y)).
\]
\end{remark}

\section{Coalgebras as curried labelled relative actions}
\label{sec:poly-relative-coalgebras-labelled-actions}

Let
\[
X\xrightarrow{j_X}J_Y=A_0\times Y
\]
be an object over the relative interface.  Write
\[
p_X:=\pi_Aj_X:X\to A_0,
\qquad
y_X:=\pi_Yj_X:X\to Y.
\]

A coalgebra
\[
c:X\to\mathcal M_{\mathbb A,\mathbb B,Y}(X)
\]
assigns to each
\[
z\in X_{(v,y)}
\]
and each
\[
a:u\to v
\]
a pair
\[
c(z)(a)=(\beta_a,z_a),
\]
where
\[
\beta_a:b'_a\to r(y)
\]
and
\[
z_a\in X_{(u,\sigma(\beta_a,y))}.
\]

Thus a coalgebra determines uncurried maps
\[
\alpha_c:
A_1\times_{t_A,A_0,p_X}X\to X
\]
and
\[
\lambda_c:
A_1\times_{t_A,A_0,p_X}X\to B_1
\]
by
\[
c(z)(a)=(\lambda_c(a,z),\alpha_c(a,z)).
\]

These satisfy the typing equations
\[
p_X\alpha_c(a,z)=s_A(a),
\]
\[
t_B\lambda_c(a,z)=r(y_Xz),
\]
and
\[
y_X\alpha_c(a,z)
=
\sigma(\lambda_c(a,z),y_Xz).
\]

Conversely, maps \((\alpha,\lambda)\) satisfying these typing equations
determine a coalgebra
\[
c:X\to\mathcal M_{\mathbb A,\mathbb B,Y}(X)
\]
by
\[
c(z)(a)=(\lambda(a,z),\alpha(a,z)).
\]

\begin{theorem}[Currying labelled relative one-step data]
\label{thm:poly-relative-currying-labelled-one-step}
To give a coalgebra
\[
c:X\to\mathcal M_{\mathbb A,\mathbb B,Y}(X)
\]
is equivalently to give maps
\[
\alpha:
A_1\times_{t_A,A_0,p_X}X\to X,
\qquad
\lambda:
A_1\times_{t_A,A_0,p_X}X\to B_1
\]
satisfying
\[
p_X\alpha=s_A\pi_1,
\]
\[
t_B\lambda=r y_X\pi_2,
\]
and
\[
y_X\alpha=\sigma(\lambda,y_X\pi_2).
\]
\end{theorem}

\begin{proof}
This is the fibre calculation of
Proposition~\ref{prop:poly-relative-fibre-calculation} read in both
directions.  A coalgebra chooses, for each input arrow \(a\), an output arrow
and a continuation state over the interface obtained by executing that output
in \(Y\).  These are precisely the maps \(\lambda\) and \(\alpha\) satisfying
the displayed typing equations.
\end{proof}

\section{Multiplicative relative response coalgebras}
\label{sec:poly-relative-multiplicative-coalgebras}

The coalgebra map records the typed one-step labelled response profile.  To
correspond to relative program categories, this profile must respect identities
and composition in the internal categories \(\mathbb A\) and \(\mathbb B\), and
the action law of \(Y\).

\begin{definition}[Multiplicative relative response coalgebra]
\label{def:poly-relative-multiplicative-response-coalgebra}
A coalgebra
\[
c:X\to\mathcal M_{\mathbb A,\mathbb B,Y}(X)
\]
is called \textbf{multiplicative} if, writing
\[
c(z)(a)=(\beta_a,z_a),
\]
the following equations hold.

\begin{enumerate}
\item For every \(z\in X\),
\[
c(z)(e_A(p_Xz))
=
(e_B(r(y_Xz)),z).
\]

\item For every composable pair
\[
a:u\to v,
\qquad
b:v\to w,
\]
and every
\[
z\in X_{(w,y)},
\]
if
\[
c(z)(b)=(\beta,z')
\]
and
\[
c(z')(a)=(\alpha,z''),
\]
then
\[
c(z)(m_A(a,b))
=
(m_B(\alpha,\beta),z'').
\]
\end{enumerate}
\end{definition}

Here the order is the convention fixed earlier: \(m_A(a,b)\) denotes the
composite of
\[
a:u\to v,
\qquad
b:v\to w,
\]
and \(m_B(\alpha,\beta)\) denotes the corresponding composite output arrow.

\begin{proposition}[Multiplicative coalgebras as labelled input-discrete categories]
\label{prop:poly-relative-coalgebras-labelled-input-discrete}
Assume the relevant dependent products exist.  To give a multiplicative
coalgebra
\[
c:X\to\mathcal M_{\mathbb A,\mathbb B,Y}(X)
\]
is equivalently to give:
\begin{enumerate}
\item a right \(\mathbb A\)-action on \(X\xrightarrow{p_X}A_0\);
\item an internal functor
\[
\Lambda_X:\int_{\mathbb A}X\to\int_{\mathbb B}Y
\]
whose object map is
\[
y_X:X\to Y.
\]
\end{enumerate}
Equivalently, it is a relative labelled input-discrete category over \(Y\).
\end{proposition}

\begin{proof}
By Theorem~\ref{thm:poly-relative-currying-labelled-one-step}, a coalgebra is
equivalent to maps
\[
\alpha:
A_1\times_{t_A,A_0,p_X}X\to X
\]
and
\[
\lambda:
A_1\times_{t_A,A_0,p_X}X\to B_1
\]
with the displayed typing equations.

The first component \(\alpha\), together with the identity and composition
parts of the multiplicativity equations, gives a right \(\mathbb A\)-action on
\(X\).  The action category \(\int_{\mathbb A}X\) has arrows \((a,z)\).

The map
\[
(a,z)\longmapsto(\lambda(a,z),y_Xz)
\]
is well-defined as a map into the arrow object of \(\int_{\mathbb B}Y\) because
\[
t_B\lambda(a,z)=r(y_Xz).
\]
Its target compatibility is exactly
\[
y_X\alpha(a,z)=\sigma(\lambda(a,z),y_Xz).
\]
The unit and multiplication clauses of
Definition~\ref{def:poly-relative-multiplicative-response-coalgebra}
are exactly identity and composition preservation for this functor.  Conversely,
such a functor gives the labelled map \(\lambda\), and its functoriality gives
the multiplicativity equations.
\end{proof}

\section{The relative coalgebra induced by a Mealy morphism}
\label{sec:poly-relative-mealy-induced-coalgebra}

Let
\[
\Phi=(P,\delta,\omega):\mathbb A\rightsquigarrow\mathbb B
\]
be a Mealy morphism and let
\[
Y=(Y\xrightarrow{r}B_0,\sigma)
\]
be a right \(\mathbb B\)-action.

Recall that
\[
S_{\Phi,Y}=P\times_{B_0}Y
\]
maps to
\[
J_Y=A_0\times Y
\]
by
\[
(x,y)\longmapsto(p(x),y).
\]

\begin{definition}[Relative Mealy response coalgebra]
\label{def:poly-relative-mealy-response-coalgebra}
The \textbf{relative Mealy response coalgebra} induced by \(\Phi\) at \(Y\) is
the coalgebra
\[
c_{\Phi,Y}:
S_{\Phi,Y}\to
\mathcal M_{\mathbb A,\mathbb B,Y}(S_{\Phi,Y})
\]
given fibrewise by
\[
c_{\Phi,Y}(x,y)(a)
=
\left(
\omega(a,x),
\left(
\delta(a,x),
\sigma(\omega(a,x),y)
\right)
\right).
\]
\end{definition}

\begin{theorem}[Relative Mealy coalgebra as transpose of labelled EM action]
\label{thm:poly-relative-mealy-coalgebra-transpose}
The coalgebra
\[
c_{\Phi,Y}
\]
is the dependent-product transpose of the labelled relative EM action
\[
(\Phi^\ast Y,\Lambda_{\Phi,Y}).
\]
Moreover, it is multiplicative.
\end{theorem}

\begin{proof}
The action map of \(\Phi^\ast Y\) is
\[
(a,x,y)
\longmapsto
\left(
\delta(a,x),
\sigma(\omega(a,x),y)
\right).
\]
The downstream comparison functor supplies the output label
\[
(a,x,y)\longmapsto \omega(a,x).
\]
Together these give the uncurried labelled one-step data
\[
(a,x,y)
\longmapsto
\left(
\omega(a,x),
\left(
\delta(a,x),
\sigma(\omega(a,x),y)
\right)
\right).
\]
By Theorem~\ref{thm:poly-relative-currying-labelled-one-step}, this transposes
to the displayed coalgebra \(c_{\Phi,Y}\).

Multiplicativity follows from the Mealy unit and multiplication equations
together with the unit and multiplication laws of the right
\(\mathbb B\)-action \(Y\).  Equivalently, it follows from
Proposition~\ref{prop:poly-relative-coalgebras-labelled-input-discrete},
since the previous chapter showed that
\[
\Lambda_{\Phi,Y}:\mathsf{Prog}_{\Phi,Y}\to\int_{\mathbb B}Y
\]
is an internal functor.
\end{proof}

\section{Position objects and flattening}
\label{sec:poly-relative-position-objects-flattening}

The polynomial coalgebra is the curried form of the relative one-step
semantics.  Flattening recovers the execution span of the relative program
category.

Let
\[
J_Y
\xleftarrow{s}
E_{\mathbb A,\mathbb B,Y}
\xrightarrow{p}
K_{\mathbb A,\mathbb B,Y}
\xrightarrow{t}
J_Y
\]
be the relative Mealy response polynomial.

For an object
\[
X\to J_Y,
\]
define the \textbf{position object of the polynomial at \(X\)} by
\[
\mathrm{Pos}_X
:=
E_{\mathbb A,\mathbb B,Y}
\times_{K_{\mathbb A,\mathbb B,Y}}
\Pi_p(s^\ast X).
\]
It carries a canonical projection
\[
\rho_X:\mathrm{Pos}_X\to
\Sigma_t\Pi_p(s^\ast X)
=
\mathcal M_{\mathbb A,\mathbb B,Y}(X).
\]
The projection is obtained by composing the second projection
\[
\mathrm{Pos}_X\to\Pi_p(s^\ast X)
\]
with the structure map of
\[
\Sigma_t\Pi_p(s^\ast X)
\]
over \(J_Y\).

For \(X=S_{\Phi,Y}\), write
\[
\mathrm{Pos}_{\Phi,Y}:=\mathrm{Pos}_{S_{\Phi,Y}},
\qquad
\rho:=\rho_{S_{\Phi,Y}}.
\]

The position object carries evaluation maps
\[
\operatorname{in}:\mathrm{Pos}_{\Phi,Y}\to A_1,
\]
\[
\operatorname{out}:\mathrm{Pos}_{\Phi,Y}\to B_1,
\]
\[
\operatorname{down}:\mathrm{Pos}_{\Phi,Y}\to
\left(\int_{\mathbb B}Y\right)_1,
\]
and
\[
\operatorname{next}:\mathrm{Pos}_{\Phi,Y}\to S_{\Phi,Y}.
\]
Externally, a point of \(\mathrm{Pos}_{\Phi,Y}\) is a total response profile
together with a chosen input position.  The evaluation maps extract the chosen
input arrow, the raw output arrow selected by the profile at that position, the
downstream transition, and the continuation state.

Now let
\[
c_{\Phi,Y}:
S_{\Phi,Y}\to
\mathcal M_{\mathbb A,\mathbb B,Y}(S_{\Phi,Y})
\]
be the relative Mealy response coalgebra.  Pulling back the position projection
along \(c_{\Phi,Y}\) gives
\[
\begin{tikzcd}[column sep=large,row sep=large]
M_{\Phi,Y}
  \arrow[r]
  \arrow[d, "s_{\Phi,Y}"']
&
\mathrm{Pos}_{\Phi,Y}
  \arrow[d, "\rho"]
\\
S_{\Phi,Y}
  \arrow[r, "c_{\Phi,Y}"']
&
\mathcal M_{\mathbb A,\mathbb B,Y}(S_{\Phi,Y}) .
\end{tikzcd}
\]

Under the canonical identification, this pullback object is
\[
M_{\Phi,Y}
=
A_1\times_{t_A,A_0,p\pi_P}S_{\Phi,Y}.
\]
A generalized element is a triple
\[
(a,x,y)
\]
with
\[
t_A(a)=p(x),
\qquad
q(x)=r(y).
\]

Under this identification, the evaluation maps become
\[
\operatorname{in}(a,x,y)=a,
\]
\[
\operatorname{out}(a,x,y)=\omega(a,x),
\]
\[
\operatorname{down}(a,x,y)=(\omega(a,x),y),
\]
and
\[
\operatorname{next}(a,x,y)
=
\left(
\delta(a,x),
\sigma(\omega(a,x),y)
\right).
\]

\begin{theorem}[Program category recovered by flattening]
\label{thm:poly-relative-program-category-recovered}
Flattening the relative Mealy response coalgebra
\[
c_{\Phi,Y}
\]
recovers the primitive execution span
\[
\mathbf m_{\Phi,Y}
=
\left(
S_{\Phi,Y}
\xleftarrow{s_{\Phi,Y}}
M_{\Phi,Y}
\xrightarrow{t_{\Phi,Y}}
S_{\Phi,Y}
\right)
\]
of the relative program category
\[
\mathsf{Prog}_{\Phi,Y}.
\]
Here
\[
s_{\Phi,Y}(a,x,y)=(x,y)
\]
and
\[
t_{\Phi,Y}(a,x,y)
=
\left(
\delta(a,x),
\sigma(\omega(a,x),y)
\right).
\]
Consequently, the internal category recovered from the flattened coalgebra is
exactly
\[
\mathsf{Prog}_{\Phi,Y}
=
\int_{\mathbb A}\Phi^\ast Y.
\]
\end{theorem}

\begin{proof}
The positions of the total response profile are precisely the admissible input
queries at a relative state.  Pulling the position object back along
\(c_{\Phi,Y}\) therefore gives the object of triples \((a,x,y)\) with
\(t_A(a)=p(x)\) and \(q(x)=r(y)\), namely
\[
A_1\times_{t_A,A_0,p\pi_P}S_{\Phi,Y}.
\]
The source is the current state \((x,y)\), and the target is the continuation
state given by the relative action \(\Phi^\ast Y\).  These are exactly the
source and target maps of the primitive execution span constructed in the
previous chapter.  The multiplicativity of \(c_{\Phi,Y}\) supplies the identity
and composition maps, so the recovered internal category is
\[
\int_{\mathbb A}\Phi^\ast Y.
\]
\end{proof}

\begin{remark}
\label{rem:poly-relative-not-replacement}
The polynomial coalgebra does not replace the program category.  It is the
curried total-response presentation of the same labelled one-step data whose
flattened form is the primitive execution span.
\end{remark}

\section{The terminal trace specialization}
\label{sec:poly-relative-terminal-specialization}

The terminal trace polynomial is now a special case of the relative response
polynomial.

Let
\[
Y=\mathbf 1_{\mathbb B}
=
(B_0\xrightarrow{1_{B_0}}B_0,\tau_B)
\]
be the terminal right \(\mathbb B\)-action.  Thus
\[
\tau_B(\beta,t_B\beta)=s_B\beta.
\]

Then
\[
J_Y=A_0\times B_0.
\]
The fibre formula of Proposition~\ref{prop:poly-relative-fibre-calculation}
becomes
\[
\mathcal M_{\mathbb A,\mathbb B,\mathbf 1_{\mathbb B}}(X)_{(v,y)}
\cong
\prod_{a:u\to v}
\sum_{\beta:y'\to y}
X_{(u,y')}.
\]

Moreover,
\[
S_{\Phi,\mathbf 1_{\mathbb B}}
=
P\times_{B_0}B_0
\cong
P.
\]
Under this isomorphism,
\[
c_{\Phi,\mathbf 1_{\mathbb B}}(x)(a)
=
(\omega(a,x),\delta(a,x)).
\]

\begin{corollary}[Terminal Mealy polynomial]
\label{cor:poly-relative-terminal-mealy-polynomial}
The terminal response polynomial of a Mealy morphism
\[
\Phi:\mathbb A\rightsquigarrow\mathbb B
\]
is the special case
\[
Y=\mathbf 1_{\mathbb B}
\]
of the relative response polynomial.  Its fibre is
\[
\prod_{a:u\to v}
\sum_{\beta:y'\to y}
X_{(u,y')}.
\]
The induced coalgebra is the terminal trace response profile
\[
x
\longmapsto
\left(
a
\longmapsto
(\omega(a,x),\delta(a,x))
\right).
\]
Flattening this coalgebra recovers the trace program category
\[
\mathsf{Prog}_\Phi
=
\mathsf{Prog}_{\Phi,\mathbf 1_{\mathbb B}}.
\]
\end{corollary}

\section{Guard loci in the relative setting}
\label{sec:poly-relative-guard-loci}

The relative polynomial presentation separates several natural loci of
guardedness.

Fix \(\Phi\) and \(Y\), and write
\[
S:=S_{\Phi,Y},
\qquad
M:=M_{\Phi,Y}.
\]

A guarded primitive specification may involve:
\[
\theta\hookrightarrow S
\qquad
\text{current coupled-state guard},
\]
\[
R\hookrightarrow A_1
\qquad
\text{input-arrow guard},
\]
\[
O\hookrightarrow B_1
\qquad
\text{raw output-arrow guard},
\]
\[
T\hookrightarrow \left(\int_{\mathbb B}Y\right)_1
\qquad
\text{downstream transition guard},
\]
and
\[
\varphi\hookrightarrow S
\qquad
\text{continuation-state predicate}.
\]

The terminal trace case has no independent downstream state.  Then
\[
\int_{\mathbb B}\mathbf 1_{\mathbb B}\cong\mathbb B^{\mathrm{op}},
\]
so downstream transition guards reduce to predicates on output arrows of
\(\mathbb B\).

\section{Guarded primitive relative programs}
\label{sec:poly-relative-guarded-primitive-programs}

The primitive execution object
\[
M_{\Phi,Y}
=
A_1\times_{t_A,A_0,p\pi_P}S_{\Phi,Y}
\]
has several canonical maps:
\[
\iota_A:M_{\Phi,Y}\to A_1,
\qquad
\iota_A(a,x,y)=a,
\]
\[
\iota_B:M_{\Phi,Y}\to B_1,
\qquad
\iota_B(a,x,y)=\omega(a,x),
\]
and
\[
\lambda_{\Phi,Y}:M_{\Phi,Y}\to
\left(\int_{\mathbb B}Y\right)_1,
\qquad
\lambda_{\Phi,Y}(a,x,y)=(\omega(a,x),y).
\]

Let
\[
\theta\hookrightarrow S_{\Phi,Y},
\qquad
R\hookrightarrow A_1,
\qquad
O\hookrightarrow B_1,
\qquad
T\hookrightarrow \left(\int_{\mathbb B}Y\right)_1
\]
be guards.

Pull them back to \(M_{\Phi,Y}\):
\[
\theta^\sharp:=s_{\Phi,Y}^\ast\theta,
\qquad
R^\sharp:=\iota_A^\ast R,
\qquad
O^\sharp:=\iota_B^\ast O,
\qquad
T^\sharp:=\lambda_{\Phi,Y}^\ast T.
\]
Their simultaneous restriction is
\[
(\theta;R/O/T)^\sharp
:=
\theta^\sharp\wedge R^\sharp\wedge O^\sharp\wedge T^\sharp
\hookrightarrow M_{\Phi,Y}.
\]

\begin{definition}[Guarded primitive relative program]
\label{def:poly-relative-guarded-primitive-relative-program}
The \textbf{guarded primitive relative program}
\[
\theta;R/O/T
\]
is the restricted execution span
\[
S_{\Phi,Y}
\xleftarrow{s_{\Phi,Y}i}
(\theta;R/O/T)^\sharp
\xrightarrow{t_{\Phi,Y}i}
S_{\Phi,Y},
\]
where
\[
i:(\theta;R/O/T)^\sharp\hookrightarrow M_{\Phi,Y}
\]
is the inclusion.  Its diamond modality is denoted
\[
\langle \theta;R/O/T\rangle\varphi.
\]
\end{definition}

\begin{proposition}[External reading]
\label{prop:poly-relative-guarded-primitive-external}
In \(\mathbf{Set}\),
\[
(x,y)\models \langle \theta;R/O/T\rangle\varphi
\]
holds precisely when
\[
\theta(x,y)
\]
and there exists an admissible input arrow
\[
a:u\to p(x)
\]
such that
\[
R(a),
\]
\[
O(\omega(a,x)),
\]
\[
T(\omega(a,x),y),
\]
and
\[
\left(
\delta(a,x),
\sigma(\omega(a,x),y)
\right)
\models\varphi.
\]
\end{proposition}

\begin{proof}
The restricted primitive span has apex the subobject of \(M_{\Phi,Y}\)
consisting of triples \((a,x,y)\) satisfying the four pulled-back guards.
The diamond says that there exists such a triple with source \((x,y)\) and
target satisfying \(\varphi\).  The target is
\[
\left(
\delta(a,x),
\sigma(\omega(a,x),y)
\right),
\]
which gives the displayed reading.
\end{proof}

\section{Guarded polynomial response liftings}
\label{sec:poly-relative-guarded-response-liftings}

Assume now that \(\mathcal E\) is regular.

Let
\[
c_{\Phi,Y}:
S_{\Phi,Y}\to
\mathcal M_{\mathbb A,\mathbb B,Y}(S_{\Phi,Y})
\]
be the relative Mealy response coalgebra, and let
\[
\rho:
\mathrm{Pos}_{\Phi,Y}\to
\mathcal M_{\mathbb A,\mathbb B,Y}(S_{\Phi,Y})
\]
be the position projection from
Section~\ref{sec:poly-relative-position-objects-flattening}.

The position object carries evaluation maps
\[
\operatorname{in}:\mathrm{Pos}_{\Phi,Y}\to A_1,
\]
\[
\operatorname{out}:\mathrm{Pos}_{\Phi,Y}\to B_1,
\]
\[
\operatorname{down}:\mathrm{Pos}_{\Phi,Y}\to
\left(\int_{\mathbb B}Y\right)_1,
\]
and
\[
\operatorname{next}:\mathrm{Pos}_{\Phi,Y}\to S_{\Phi,Y}.
\]

Externally, a point of \(\mathrm{Pos}_{\Phi,Y}\) is a total response profile
together with a chosen input position.  The evaluation maps extract:
\[
\operatorname{in}=\text{the chosen input arrow},
\]
\[
\operatorname{out}=\text{the raw output arrow},
\]
\[
\operatorname{down}=\text{the downstream transition},
\]
and
\[
\operatorname{next}=\text{the continuation state}.
\]

\begin{definition}[Guarded existential response lifting]
\label{def:poly-relative-guarded-existential-lifting}
Let
\[
R\hookrightarrow A_1,
\qquad
O\hookrightarrow B_1,
\qquad
T\hookrightarrow\left(\int_{\mathbb B}Y\right)_1,
\qquad
\varphi\hookrightarrow S_{\Phi,Y}
\]
be predicates.  Define the guarded existential response lifting
\[
\lambda_{R/O/T}(\varphi)
\hookrightarrow
\mathcal M_{\mathbb A,\mathbb B,Y}(S_{\Phi,Y})
\]
by
\[
\lambda_{R/O/T}(\varphi)
:=
\exists_{\rho}
\left(
\operatorname{in}^{\ast}R
\wedge
\operatorname{out}^{\ast}O
\wedge
\operatorname{down}^{\ast}T
\wedge
\operatorname{next}^{\ast}\varphi
\right).
\]
\end{definition}

Externally, a response profile satisfies
\[
\lambda_{R/O/T}(\varphi)
\]
when some input position satisfies \(R\), its raw output satisfies \(O\), its
downstream transition satisfies \(T\), and its continuation state satisfies
\(\varphi\).

\begin{theorem}[Agreement of polynomial liftings and span modalities]
\label{thm:poly-relative-lifting-span-agreement}
For every state guard
\[
\theta\hookrightarrow S_{\Phi,Y},
\]
input guard
\[
R\hookrightarrow A_1,
\]
output guard
\[
O\hookrightarrow B_1,
\]
downstream transition guard
\[
T\hookrightarrow\left(\int_{\mathbb B}Y\right)_1,
\]
and postcondition
\[
\varphi\hookrightarrow S_{\Phi,Y},
\]
one has
\[
\theta\wedge c_{\Phi,Y}^{\ast}\lambda_{R/O/T}(\varphi)
=
\langle \theta;R/O/T\rangle\varphi.
\]
\end{theorem}

\begin{proof}
Pulling back the position projection
\[
\rho:\mathrm{Pos}_{\Phi,Y}\to
\mathcal M_{\mathbb A,\mathbb B,Y}(S_{\Phi,Y})
\]
along the coalgebra \(c_{\Phi,Y}\) gives the primitive execution object
\[
M_{\Phi,Y}.
\]
Under this identification, the evaluation maps become
\[
\operatorname{in}=\iota_A,
\qquad
\operatorname{out}=\iota_B,
\qquad
\operatorname{down}=\lambda_{\Phi,Y},
\qquad
\operatorname{next}=t_{\Phi,Y}.
\]
By Beck--Chevalley for regular images,
\[
c_{\Phi,Y}^{\ast}\lambda_{R/O/T}(\varphi)
=
\exists_{s_{\Phi,Y}}
\left(
\iota_A^\ast R
\wedge
\iota_B^\ast O
\wedge
\lambda_{\Phi,Y}^\ast T
\wedge
t_{\Phi,Y}^\ast\varphi
\right).
\]
Using Frobenius reciprocity,
\[
\theta\wedge c_{\Phi,Y}^{\ast}\lambda_{R/O/T}(\varphi)
=
\exists_{s_{\Phi,Y}}
\left(
s_{\Phi,Y}^\ast\theta
\wedge
\iota_A^\ast R
\wedge
\iota_B^\ast O
\wedge
\lambda_{\Phi,Y}^\ast T
\wedge
t_{\Phi,Y}^\ast\varphi
\right).
\]
The right-hand side is exactly the diamond of the restricted primitive span
\[
\theta;R/O/T.
\]
\end{proof}

\begin{remark}[Central comparison]
\label{rem:poly-relative-central-comparison}
Theorem~\ref{thm:poly-relative-lifting-span-agreement} is the central bridge
of the chapter.  The polynomial coalgebra gives a guarded one-step response
lifting on total response profiles.  Pulling that lifting back along the
relative Mealy coalgebra recovers the guarded span diamond of the relative
program category.
\end{remark}

\section{Profile modalities}
\label{sec:poly-relative-profile-modalities}

The existential lifting of
Definition~\ref{def:poly-relative-guarded-existential-lifting} recovers the
ordinary guarded span diamond after pulling back along the coalgebra.  The
polynomial presentation also displays modalities that quantify directly over
the input positions of a total response profile.

Assume \(\mathcal E\) is a first-order Heyting category.

Let
\[
R\hookrightarrow A_1,
\qquad
O\hookrightarrow B_1,
\qquad
T\hookrightarrow\left(\int_{\mathbb B}Y\right)_1,
\qquad
\varphi\hookrightarrow S_{\Phi,Y}
\]
be predicates.

Define the \textbf{weak universal profile modality} by
\[
\Box^{\mathrm{prof}}_{R/O/T}(\varphi)
:=
\forall_{\rho}
\left(
(\operatorname{in}^{\ast}R
\wedge
\operatorname{out}^{\ast}O
\wedge
\operatorname{down}^{\ast}T)
\Rightarrow
\operatorname{next}^{\ast}\varphi
\right).
\]
Externally, this says:
\[
\forall a:
\quad
R(a)\wedge O(\operatorname{out}(a))
\wedge T(\operatorname{down}(a))
\Rightarrow
\varphi(\operatorname{next}(a)).
\]

Define the \textbf{strong universal profile modality} by
\[
\Box^{\mathrm{prof,strong}}_{R/O/T}(\varphi)
:=
\forall_{\rho}
\left(
\operatorname{in}^{\ast}R
\Rightarrow
\left(
\operatorname{out}^{\ast}O
\wedge
\operatorname{down}^{\ast}T
\wedge
\operatorname{next}^{\ast}\varphi
\right)
\right).
\]
Externally, this says:
\[
\forall a:
\quad
R(a)
\Rightarrow
\bigl(
O(\operatorname{out}(a))
\wedge
T(\operatorname{down}(a))
\wedge
\varphi(\operatorname{next}(a))
\bigr).
\]

\begin{theorem}[Agreement of weak profile boxes and span boxes]
\label{thm:poly-relative-profile-box-span-agreement}
Assume \(\mathcal E\) is a first-order Heyting category.  Then
\[
c_{\Phi,Y}^{\ast}
\Box^{\mathrm{prof}}_{R/O/T}(\varphi)
=
[R/O/T]\varphi,
\]
where \([R/O/T]\varphi\) denotes the box modality of the restricted primitive
span determined by the guards \(R,O,T\).

More generally, for a state guard
\[
\theta\hookrightarrow S_{\Phi,Y},
\]
one has
\[
\theta\Rightarrow
c_{\Phi,Y}^{\ast}
\Box^{\mathrm{prof}}_{R/O/T}(\varphi)
=
[\theta;R/O/T]\varphi.
\]
\end{theorem}

\begin{proof}
By Beck--Chevalley for the right adjoints \(\forall\), pulling back
\(\Box^{\mathrm{prof}}_{R/O/T}(\varphi)\) along \(c_{\Phi,Y}\) gives
\[
\forall_{s_{\Phi,Y}}
\left(
(\iota_A^\ast R
\wedge
\iota_B^\ast O
\wedge
\lambda_{\Phi,Y}^\ast T)
\Rightarrow
t_{\Phi,Y}^\ast\varphi
\right).
\]
This is exactly the box of the restricted primitive span determined by
\(R,O,T\).

For the guarded version, let
\[
U:=(R/O/T)^\sharp\hookrightarrow M_{\Phi,Y}.
\]
Then
\[
[\theta;R/O/T]\varphi
=
[\theta?;U]\varphi
=
\theta\Rightarrow[U]\varphi
=
\theta\Rightarrow
c_{\Phi,Y}^{\ast}
\Box^{\mathrm{prof}}_{R/O/T}(\varphi),
\]
using the test-prefix box law.
\end{proof}

\begin{remark}[Profile modalities versus span boxes]
\label{rem:poly-relative-profile-versus-span-boxes}
The weak universal profile modality is the profile-level predicate whose
pullback along the coalgebra recovers the span box of the corresponding
restricted primitive span.

The strong profile modality is different.  It requires every \(R\)-input
position to produce an output satisfying \(O\), a downstream transition
satisfying \(T\), and a continuation satisfying \(\varphi\).  Thus the strong
profile modality is generally stricter than the box of the restricted span.
\end{remark}

\section{Tests, guarded commands, and guarded choice}
\label{sec:poly-relative-tests-guarded-choice}

The remaining guarded program operations are inherited from the span dynamic
logic of Chapter~\ref{ch:relative-mealy-program-categories}.  We record them
for the relative state object
\[
S:=S_{\Phi,Y}.
\]

Assume \(\mathcal E\) is regular.  For a predicate
\[
\theta:\Theta\hookrightarrow S,
\]
the corresponding test span is
\[
\theta?
=
\left(
S\xleftarrow{\theta}\Theta\xrightarrow{\theta}S
\right).
\]

Let
\[
U:S\nrightarrow S
\]
be an endoprogram.  The guarded command
\[
\theta?;U
\]
is the sequential composite of the test and \(U\).

\begin{proposition}[Guard-prefix law]
\label{prop:poly-relative-guard-prefix-law}
For every endoprogram \(U:S\nrightarrow S\),
\[
\langle \theta?;U\rangle\varphi
=
\theta\wedge\langle U\rangle\varphi.
\]
\end{proposition}

\begin{proof}
By functoriality of diamonds,
\[
\langle \theta?;U\rangle\varphi
=
\langle \theta?\rangle(\langle U\rangle\varphi).
\]
The test law gives
\[
\langle \theta?\rangle\psi=\theta\wedge\psi.
\]
Substituting \(\psi=\langle U\rangle\varphi\) gives the result.
\end{proof}

Assume now that \(\mathcal E\) is coherent.  Let
\[
T:S\nrightarrow S
\]
be a fixed ambient endospan, and let
\[
U,V\hookrightarrow T
\]
be subprograms.  If
\[
\theta,\epsilon\hookrightarrow S
\]
are state guards, then
\[
\theta?;U
\qquad
\text{and}
\qquad
\epsilon?;V
\]
are identified, under the canonical pullback isomorphisms for composing with
test spans, with the subobjects
\[
\ell_U^\ast\theta\wedge U\hookrightarrow T
\qquad
\text{and}
\qquad
\ell_V^\ast\epsilon\wedge V\hookrightarrow T,
\]
where \(\ell_U\) and \(\ell_V\) are the source legs of \(U\) and \(V\).  Their
guarded choice is the join
\[
(\theta?;U)\vee(\epsilon?;V)\hookrightarrow T.
\]

\begin{proposition}[Diamond law for guarded choice]
\label{prop:poly-relative-guarded-choice-diamond}
For guarded branches inside a common ambient endospan,
\[
\left\langle
(\theta?;U)\vee(\epsilon?;V)
\right\rangle\varphi
=
(\theta\wedge\langle U\rangle\varphi)
\vee
(\epsilon\wedge\langle V\rangle\varphi).
\]
\end{proposition}

\begin{proof}
This is the finite choice law for subprograms in a coherent category, followed
by the guard-prefix law.
\end{proof}

\begin{definition}[Internal conditional]
\label{def:poly-relative-internal-conditional}
Let
\[
T:S\nrightarrow S
\]
be a fixed ambient endospan, and let
\[
U,V\hookrightarrow T
\]
be subprograms.  Let
\[
\theta,\epsilon\hookrightarrow S
\]
be complementary guards.  Define
\[
\mathbf{if}\ \theta\ \mathbf{then}\ U\ \mathbf{else}\ V
:=
(\theta?;U)\vee(\epsilon?;V)
\hookrightarrow T.
\]
In Boolean fibres, one may take
\[
\epsilon=\neg\theta.
\]
\end{definition}

\section{Boxes for guarded choice and conditionals}
\label{sec:poly-relative-boxes-guarded-choice}

Assume \(\mathcal E\) is a first-order Heyting category.

For a test,
\[
[\theta?]\varphi=\theta\Rightarrow\varphi.
\]
Therefore
\[
[\theta?;U]\varphi
=
\theta\Rightarrow[U]\varphi.
\]

\begin{proposition}[Box of finite subprogram choice]
\label{prop:poly-relative-box-finite-choice}
For subprograms
\[
W,Z\hookrightarrow T
\]
inside a fixed ambient endospan \(T:S\nrightarrow S\),
\[
[W\vee Z]\varphi
=
[W]\varphi\wedge[Z]\varphi.
\]
\end{proposition}

\begin{proof}
Regard \(W\) and \(Z\) as predicates on the apex of \(T\).  If
\[
T=(S\xleftarrow{\ell_T}T\xrightarrow{r_T}S),
\]
then the box of a restricted subspan \(W\hookrightarrow T\) may be computed as
\[
[W]\varphi
=
\forall_{\ell_T}\bigl(W\Rightarrow r_T^\ast\varphi\bigr).
\]
Indeed, this formula is equivalent, by adjunction, to requiring
\[
\ell_T^\ast\psi\wedge W\le r_T^\ast\varphi
\]
for every predicate \(\psi\hookrightarrow S\), which is precisely the defining
condition for the restricted span.

Hence
\[
[W\vee Z]\varphi
=
\forall_{\ell_T}\bigl((W\vee Z)\Rightarrow r_T^\ast\varphi\bigr).
\]
In a Heyting algebra,
\[
(W\vee Z)\Rightarrow A
=
(W\Rightarrow A)\wedge(Z\Rightarrow A).
\]
Since \(\forall_{\ell_T}\) is a right adjoint, it preserves finite meets.
\end{proof}

\begin{proposition}[Box law for conditionals]
\label{prop:poly-relative-conditional-box-law}
For complementary guards \(\theta,\epsilon\) and subprograms
\[
U,V\hookrightarrow T
\]
inside a common ambient endospan,
\[
\left[
\mathbf{if}\ \theta\ \mathbf{then}\ U\ \mathbf{else}\ V
\right]\varphi
=
(\theta\Rightarrow[U]\varphi)
\wedge
(\epsilon\Rightarrow[V]\varphi).
\]
\end{proposition}

\begin{proof}
The conditional is
\[
(\theta?;U)\vee(\epsilon?;V).
\]
Apply Proposition~\ref{prop:poly-relative-box-finite-choice} and the test-prefix
box law.
\end{proof}

\section{Guarded iteration and while programs}
\label{sec:poly-relative-while-programs}

Assume now that \(\mathcal E\) has the path-star structure and fixed-point
semantics of Chapter~\ref{ch:relative-mealy-program-categories}.  Let
\[
U:S\nrightarrow S
\]
be an endoprogram, and let
\[
\theta,\epsilon\hookrightarrow S
\]
be guards.

\begin{definition}[Guarded while program]
\label{def:poly-relative-guarded-while-program}
The guarded while program with loop guard \(\theta\), body \(U\), and exit
guard \(\epsilon\) is
\[
\mathbf{while}\ \theta\ \mathbf{do}\ U\ \mathbf{exit}\ \epsilon
:=
(\theta?;U)^\star_{\mathrm{path}};\epsilon?.
\]
In Boolean predicate fibres, the ordinary notation
\[
\mathbf{while}\ \theta\ \mathbf{do}\ U
\]
denotes the special case
\[
(\theta?;U)^\star_{\mathrm{path}};(\neg\theta)?.
\]
\end{definition}

\begin{theorem}[Diamond law for guarded while programs]
\label{thm:poly-relative-while-diamond}
In an \(\omega\)-iterative first-order Heyting category,
\[
\left\langle
(\theta?;U)^\star_{\mathrm{path}};\epsilon?
\right\rangle\varphi
=
\mu X.\bigl(
\epsilon\wedge\varphi
\vee
\theta\wedge\langle U\rangle X
\bigr).
\]
\end{theorem}

\begin{proof}
By functoriality of diamonds,
\[
\left\langle
(\theta?;U)^\star_{\mathrm{path}};\epsilon?
\right\rangle\varphi
=
\left\langle
(\theta?;U)^\star_{\mathrm{path}}
\right\rangle
(\epsilon\wedge\varphi).
\]
The path/fixed-point comparison gives
\[
\left\langle
(\theta?;U)^\star_{\mathrm{path}}
\right\rangle
(\epsilon\wedge\varphi)
=
\mu X.\bigl(
\epsilon\wedge\varphi
\vee
\langle\theta?;U\rangle X
\bigr).
\]
Using
\[
\langle\theta?;U\rangle X
=
\theta\wedge\langle U\rangle X
\]
gives the displayed equation.
\end{proof}

\begin{theorem}[Box law for guarded while programs]
\label{thm:poly-relative-while-box}
In an \(\omega\)-iterative first-order Heyting category,
\[
\left[
(\theta?;U)^\star_{\mathrm{path}};\epsilon?
\right]\varphi
=
\nu X.\bigl(
(\epsilon\Rightarrow\varphi)
\wedge
(\theta\Rightarrow[U]X)
\bigr).
\]
\end{theorem}

\begin{proof}
The proof is dual to that of
Theorem~\ref{thm:poly-relative-while-diamond}, using box functoriality, the
test law
\[
[\epsilon?]\varphi=\epsilon\Rightarrow\varphi,
\]
and the path/fixed-point comparison for boxes.
\end{proof}

\section{Domain and dynamic tests}
\label{sec:poly-relative-domain-tests}

Let
\[
U:S\nrightarrow S
\]
be an endoprogram.  Its domain predicate is
\[
\operatorname{dom}(U):=\langle U\rangle\top.
\]
Thus \(\operatorname{dom}(U)\) is the predicate of states from which some
\(U\)-execution is possible.

\begin{proposition}[Domain-prefix law]
\label{prop:poly-relative-domain-prefix}
For every endoprogram \(U\),
\[
\langle \operatorname{dom}(U)?;U\rangle\varphi
=
\langle U\rangle\varphi.
\]
\end{proposition}

\begin{proof}
By the guard-prefix law,
\[
\langle \operatorname{dom}(U)?;U\rangle\varphi
=
\operatorname{dom}(U)\wedge\langle U\rangle\varphi.
\]
Since
\[
\langle U\rangle\varphi\le \langle U\rangle\top
=
\operatorname{dom}(U),
\]
the meet is \(\langle U\rangle\varphi\).
\end{proof}

If the predicate fibres are Boolean, define the antidomain predicate by
\[
\operatorname{adom}(U):=\neg\operatorname{dom}(U).
\]

\begin{proposition}[Antidomain annihilation]
\label{prop:poly-relative-antidomain-annihilation}
In Boolean predicate fibres,
\[
\langle \operatorname{adom}(U)?;U\rangle\varphi
=
\bot.
\]
\end{proposition}

\begin{proof}
By the guard-prefix law,
\[
\langle \operatorname{adom}(U)?;U\rangle\varphi
=
\operatorname{adom}(U)\wedge\langle U\rangle\varphi.
\]
Since
\[
\langle U\rangle\varphi\le\operatorname{dom}(U)
\]
and
\[
\operatorname{adom}(U)=\neg\operatorname{dom}(U),
\]
the meet is bottom.
\end{proof}

\section{Relative guarded Mealy automata}
\label{sec:poly-relative-guarded-mealy-automata}

We now isolate the finite guarded one-step structure generated by a relative
Mealy system.

Let
\[
\Phi:\mathbb A\rightsquigarrow\mathbb B
\]
be a Mealy morphism and let \(Y\) be a right \(\mathbb B\)-action.  Let
\[
(\theta_i)_{i\in I}
\]
be a finite family of state guards on \(S_{\Phi,Y}\),
\[
(R_i)_{i\in I}
\]
a finite family of input-arrow guards on \(A_1\),
\[
(O_i)_{i\in I}
\]
a finite family of output-arrow guards on \(B_1\), and
\[
(T_i)_{i\in I}
\]
a finite family of downstream transition guards on
\[
\left(\int_{\mathbb B}Y\right)_1.
\]

Each quadruple determines a guarded primitive subprogram
\[
G_i:=(\theta_i;R_i/O_i/T_i)^\sharp\hookrightarrow M_{\Phi,Y}.
\]

\begin{definition}[Relative guarded Mealy automaton]
\label{def:poly-relative-guarded-mealy-automaton}
A \textbf{relative guarded Mealy automaton} over \((\Phi,Y)\) is a finite
family of guarded primitive subprograms
\[
G_i=(\theta_i;R_i/O_i/T_i)^\sharp\hookrightarrow M_{\Phi,Y}.
\]
Its associated one-step program is the finite join
\[
G:=\bigvee_{i\in I}G_i\hookrightarrow M_{\Phi,Y}.
\]
\end{definition}

\begin{proposition}[Diamond semantics of relative guarded Mealy automata]
\label{prop:poly-relative-guarded-automaton-diamond}
For a relative guarded Mealy automaton with one-step program
\[
G=\bigvee_{i\in I}G_i,
\]
one has
\[
\langle G\rangle\varphi
=
\bigvee_{i\in I}\langle G_i\rangle\varphi.
\]
\end{proposition}

\begin{proof}
This is the finite choice law for subprograms inside the common ambient object
\(M_{\Phi,Y}\).
\end{proof}

A relative guarded Mealy automaton is \textbf{witness-disjoint} if
\[
G_i\wedge G_j=\bot
\qquad
(i\ne j).
\]
It is \textbf{branch-disjoint} if
\[
\operatorname{dom}(G_i)\wedge\operatorname{dom}(G_j)=\bot
\qquad
(i\ne j).
\]
It is \textbf{guard-total} if
\[
\bigvee_{i\in I}\operatorname{dom}(G_i)=\top.
\]
When branch-disjointness and guard-totality hold, the branch domains form a
finite disjoint cover of the state object.  Thus the enabled branch is unique
at the level of branch domains, though a chosen branch may still contain
several execution witnesses.

\section{Output, downstream, and interface guarded loops}
\label{sec:poly-relative-output-downstream-interface-guards}

The relative Mealy setting allows guards to inspect emitted output arrows and
the downstream transitions caused by those outputs.

Let
\[
U\hookrightarrow M_{\Phi,Y}
\]
be a primitive subprogram.

For an output guard
\[
O\hookrightarrow B_1,
\]
define the output-restricted body
\[
U_O:=U\wedge\iota_B^\ast O\hookrightarrow M_{\Phi,Y}.
\]

For a downstream transition guard
\[
T\hookrightarrow\left(\int_{\mathbb B}Y\right)_1,
\]
define the downstream-restricted body
\[
U_T:=U\wedge\lambda_{\Phi,Y}^\ast T\hookrightarrow M_{\Phi,Y}.
\]

For loop and exit guards
\[
\theta,\epsilon\hookrightarrow S_{\Phi,Y},
\]
one may form output-guarded and downstream-guarded while programs:
\[
(\theta?;U_O)^\star_{\mathrm{path}};\epsilon?,
\]
and
\[
(\theta?;U_T)^\star_{\mathrm{path}};\epsilon?.
\]

Their diamond laws are the corresponding instances of
Theorem~\ref{thm:poly-relative-while-diamond}:
\[
\left\langle
(\theta?;U_O)^\star_{\mathrm{path}};\epsilon?
\right\rangle\varphi
=
\mu X.\bigl(
\epsilon\wedge\varphi
\vee
\theta\wedge\langle U_O\rangle X
\bigr),
\]
and
\[
\left\langle
(\theta?;U_T)^\star_{\mathrm{path}};\epsilon?
\right\rangle\varphi
=
\mu X.\bigl(
\epsilon\wedge\varphi
\vee
\theta\wedge\langle U_T\rangle X
\bigr).
\]

Guards may also be placed directly on relative interfaces.  Let
\[
\Gamma\hookrightarrow J_Y=A_0\times Y
\]
be an interface guard.  Pulling it back along
\[
j_{\Phi,Y}:S_{\Phi,Y}\to J_Y
\]
gives a state guard
\[
j_{\Phi,Y}^\ast\Gamma\hookrightarrow S_{\Phi,Y}.
\]
Thus one may form interface-guarded primitives such as
\[
(j_{\Phi,Y}^\ast\Gamma;R/O/T)^\sharp
\hookrightarrow M_{\Phi,Y}.
\]

\begin{proposition}[External reading of interface-guarded primitives]
\label{prop:poly-relative-interface-guarded-external}
In \(\mathbf{Set}\),
\[
(x,y)\models
\langle j_{\Phi,Y}^\ast\Gamma;R/O/T\rangle\varphi
\]
holds precisely when
\[
\Gamma(p(x),y)
\]
and there exists an admissible input arrow
\[
a:u\to p(x)
\]
such that
\[
R(a),
\qquad
O(\omega(a,x)),
\qquad
T(\omega(a,x),y),
\]
and
\[
\left(
\delta(a,x),
\sigma(\omega(a,x),y)
\right)
\models\varphi.
\]
\end{proposition}

\begin{proof}
This is Proposition~\ref{prop:poly-relative-guarded-primitive-external} with
\[
\theta=j_{\Phi,Y}^\ast\Gamma.
\]
\end{proof}

\begin{remark}[Terminal interface guards]
\label{rem:poly-relative-terminal-interface-guards}
In the terminal case
\[
Y=\mathbf 1_{\mathbb B}=B_0,
\]
the interface object
\[
J_Y=A_0\times B_0
\]
is the object of input-output interface pairs.  Thus the interface-guarded
programs above specialize to the terminal interface guards on
\[
A_0\times B_0.
\]
\end{remark}

\section{The guarded program algebra generated by a relative Mealy system}
\label{sec:poly-relative-guarded-program-algebra}

The constructions above generate a guarded program calculus over the relative
state object
\[
S_{\Phi,Y}.
\]

\begin{definition}[Relative guarded Mealy program algebra]
\label{def:poly-relative-guarded-program-algebra}
The \textbf{relative guarded Mealy program algebra} generated by
\[
(\Phi,Y)
\]
is the class of endoprograms on
\[
S_{\Phi,Y}
\]
generated from:
\begin{enumerate}
\item tests \(\theta?\) for predicates
\[
\theta\hookrightarrow S_{\Phi,Y};
\]
\item guarded primitive programs
\[
\theta;R/O/T;
\]
\item sequential composition of spans;
\item finite guarded joins of subprograms inside a specified common ambient
span;
\item path-star guarded iteration
\[
(\theta?;U)^\star_{\mathrm{path}};\epsilon?.
\]
\end{enumerate}
\end{definition}

\begin{remark}[Common ambient spans]
\label{rem:poly-relative-common-ambient-spans}
Finite guarded joins in this definition are joins of subobjects inside a fixed
ambient execution span.  Choice between arbitrary generated spans requires
additional coproduct structure and should be interpreted as coproduct of span
apexes.  That is a different, witnessful operation.
\end{remark}

\begin{remark}
\label{rem:poly-relative-program-algebra-from-polynomial}
The polynomial response profile supplies the curried one-step account of
guarded primitive programs.  Once flattened, the span calculus supplies the
usual dynamic operations: tests, sequential composition, finite subprogram
choice, path iteration, diamonds, boxes, and fixed points.
\end{remark}

\section{Representable and stateless relative coalgebras}
\label{sec:poly-relative-representable-stateless}

Let
\[
F:\mathbb A\to\mathbb B
\]
be an internal functor with object map
\[
F_0:A_0\to B_0
\]
and arrow map
\[
F_1:A_1\to B_1.
\]
As in the previous chapters, \(F\) determines a representable Mealy morphism
whose carrier span is
\[
A_0\xleftarrow{1_{A_0}}A_0\xrightarrow{F_0}B_0.
\]

Let
\[
Y=(Y\xrightarrow{r}B_0,\sigma)
\]
be a right \(\mathbb B\)-action.  The relative state object is
\[
S_{F,Y}:=A_0\times_{F_0,B_0,r}Y.
\]
A generalized element is a pair
\[
(v,y)
\]
with
\[
F_0(v)=r(y).
\]

The induced relative coalgebra
\[
c_{F,Y}:
S_{F,Y}\to
\mathcal M_{\mathbb A,\mathbb B,Y}(S_{F,Y})
\]
is given by
\[
c_{F,Y}(v,y)(a)
=
\left(
F_1(a),
\left(
s_A(a),
\sigma(F_1(a),y)
\right)
\right),
\]
for every arrow
\[
a:u\to v
\]
of \(\mathbb A\).

\begin{proposition}[Stateless relative response coalgebra]
\label{prop:poly-relative-stateless-response-coalgebra}
The coalgebra
\[
c_{F,Y}
\]
is the relative Mealy response coalgebra associated to the representable Mealy
morphism induced by \(F\).  Its flattened action category is the action
category of the ordinary restricted action
\[
F^\ast Y.
\]
\end{proposition}

\begin{proof}
For the representable carrier, the state component is forced:
\[
\delta(a,v)=s_A(a).
\]
The emitted output is
\[
\omega(a,v)=F_1(a).
\]
Thus the relative continuation state is
\[
(s_A(a),\sigma(F_1(a),y)).
\]
This gives exactly the displayed coalgebra.  Flattening recovers the restricted
right \(\mathbb A\)-action along the internal functor \(F\).
\end{proof}

\begin{corollary}[Terminal stateless coalgebra]
\label{cor:poly-relative-terminal-stateless-coalgebra}
For
\[
Y=\mathbf 1_{\mathbb B},
\]
one has
\[
S_{F,\mathbf 1_{\mathbb B}}\cong A_0,
\]
and the coalgebra becomes
\[
c_F(v)(a)=(F_1(a),s_A(a)).
\]
The recovered program category is canonically
\[
\mathbb A^{\mathrm{op}},
\]
and the output labelling functor is
\[
F^{\mathrm{op}}:\mathbb A^{\mathrm{op}}\to\mathbb B^{\mathrm{op}}.
\]
\end{corollary}

\begin{proof}
In the terminal case, the relative state object is
\[
A_0\times_{F_0,B_0,1_{B_0}}B_0\cong A_0.
\]
The terminal action sends
\[
F_1(a):F_0(s_Aa)\to F_0(t_Aa)
\]
to its source object \(F_0(s_Aa)\), so the continuation state is
\(s_A(a)\).  The program category is the terminal stateless trace category,
namely \(\mathbb A^{\mathrm{op}}\), with output labelling \(F^{\mathrm{op}}\).
\end{proof}

\section{Progressive consolidation}
\label{sec:poly-relative-progressive-consolidation}

\begin{theorem}[Polynomial transposes of relative Mealy actions]
\label{thm:poly-relative-main-consolidation}
The constructions of this chapter form the following hierarchy.

\begin{enumerate}
\item A right \(\mathbb A\)-action
\[
(X\xrightarrow{p}A_0,\alpha)
\]
is an Eilenberg--Moore algebra for the span monad associated to
\(\mathbb A\).

\item When the relevant dependent product exists, such an action admits a
polynomial transpose
\[
X\to\mathcal R_{\mathbb A}(X),
\]
where
\[
\mathcal R_{\mathbb A}(X)_v
\cong
\prod_{a:u\to v}X_u.
\]

\item For a right \(\mathbb B\)-action \(Y\), the relative interface object is
\[
J_Y=A_0\times Y.
\]

\item The triple
\[
(\mathbb A,\mathbb B,Y)
\]
determines a relative Mealy response polynomial
\[
\mathcal M_{\mathbb A,\mathbb B,Y}:
\mathcal E/J_Y\to\mathcal E/J_Y
\]
with fibre
\[
\mathcal M_{\mathbb A,\mathbb B,Y}(X)_{(v,y)}
\cong
\prod_{a:u\to v}
\sum_{\beta:b'\to r(y)}
X_{(u,\sigma(\beta,y))}.
\]

\item Coalgebras
\[
X\to\mathcal M_{\mathbb A,\mathbb B,Y}(X)
\]
are equivalently labelled one-step relative action data.

\item Multiplicative coalgebras are equivalently relative labelled
input-discrete categories over \(Y\), namely right \(\mathbb A\)-actions \(X\)
equipped with internal functors
\[
\int_{\mathbb A}X\to\int_{\mathbb B}Y.
\]

\item A Mealy morphism
\[
\Phi:\mathbb A\rightsquigarrow\mathbb B
\]
and a right \(\mathbb B\)-action \(Y\) induce a canonical multiplicative
coalgebra
\[
c_{\Phi,Y}:
S_{\Phi,Y}\to
\mathcal M_{\mathbb A,\mathbb B,Y}(S_{\Phi,Y})
\]
defined by
\[
c_{\Phi,Y}(x,y)(a)
=
\left(
\omega(a,x),
\left(
\delta(a,x),
\sigma(\omega(a,x),y)
\right)
\right).
\]

\item This coalgebra is the dependent-product transpose of the labelled
relative Eilenberg--Moore action
\[
(\Phi^\ast Y,\Lambda_{\Phi,Y}).
\]

\item Flattening
\[
c_{\Phi,Y}
\]
recovers the primitive execution span
\[
\mathbf m_{\Phi,Y}
\]
and hence the relative program category
\[
\mathsf{Prog}_{\Phi,Y}
=
\int_{\mathbb A}\Phi^\ast Y.
\]

\item If \(\mathcal E\) is regular, guarded polynomial response liftings pull
back along \(c_{\Phi,Y}\) to guarded span diamonds:
\[
\theta\wedge c_{\Phi,Y}^{\ast}\lambda_{R/O/T}(\varphi)
=
\langle \theta;R/O/T\rangle\varphi.
\]

\item If \(\mathcal E\) is first-order Heyting, the weak universal profile box
pulls back along \(c_{\Phi,Y}\) to the corresponding span box:
\[
c_{\Phi,Y}^{\ast}
\Box^{\mathrm{prof}}_{R/O/T}(\varphi)
=
[R/O/T]\varphi.
\]
The strong profile box is generally stricter.

\item Tests, guarded commands, finite guarded choice, while programs, domain
operations, and guarded automata are inherited from the span dynamic logic on
the relative state object
\[
S_{\Phi,Y}.
\]

\item The terminal trace polynomial is recovered by taking
\[
Y=\mathbf 1_{\mathbb B}.
\]
Its fibre is
\[
\prod_{a:u\to v}
\sum_{\beta:y'\to y}
X_{(u,y')},
\]
and its coalgebra is
\[
x\longmapsto
\left(
a\longmapsto(\omega(a,x),\delta(a,x))
\right).
\]

\item Representable Mealy morphisms induced by internal functors
\[
F:\mathbb A\to\mathbb B
\]
give stateless relative coalgebras
\[
c_{F,Y}(v,y)(a)
=
\left(
F_1(a),
\left(
s_A(a),
\sigma(F_1(a),y)
\right)
\right).
\]
In the terminal case this reduces to
\[
c_F(v)(a)=(F_1(a),s_A(a)),
\]
and the recovered program category is
\[
\mathbb A^{\mathrm{op}}.
\]
\end{enumerate}
\end{theorem}

\begin{proof}
Part (1) is the Eilenberg--Moore description of right internal actions.
Part (2) is Proposition~\ref{prop:poly-relative-actions-as-response-coalgebras}.
Parts (3) and (4) are the construction of the relative interface and
Proposition~\ref{prop:poly-relative-fibre-calculation}.  Part (5) is
Theorem~\ref{thm:poly-relative-currying-labelled-one-step}.  Part (6) is
Proposition~\ref{prop:poly-relative-coalgebras-labelled-input-discrete}.  Parts
(7) and (8) are Definition~\ref{def:poly-relative-mealy-response-coalgebra}
and Theorem~\ref{thm:poly-relative-mealy-coalgebra-transpose}.  Part (9) is
Theorem~\ref{thm:poly-relative-program-category-recovered}.  Part (10) is
Theorem~\ref{thm:poly-relative-lifting-span-agreement}.  Part (11) is
Theorem~\ref{thm:poly-relative-profile-box-span-agreement} and
Remark~\ref{rem:poly-relative-profile-versus-span-boxes}.  Part (12) is the
span dynamic logic recalled in
Sections~\ref{sec:poly-relative-tests-guarded-choice}--\ref{sec:poly-relative-guarded-program-algebra}.
Part (13) is Corollary~\ref{cor:poly-relative-terminal-mealy-polynomial}.
Part (14) is Proposition~\ref{prop:poly-relative-stateless-response-coalgebra}
and Corollary~\ref{cor:poly-relative-terminal-stateless-coalgebra}.
\end{proof}

\begin{remark}[Final summary]
\label{rem:poly-relative-final-summary}
The polynomial structure in this chapter is not imposed from outside the
relative span semantics.  It is obtained by currying the labelled
Eilenberg--Moore action associated to a Mealy morphism and a downstream action:
\[
(\Phi^\ast Y,\Lambda_{\Phi,Y})
\quad
\leadsto
\quad
c_{\Phi,Y}.
\]
The resulting coalgebra
\[
c_{\Phi,Y}:
S_{\Phi,Y}\to
\mathcal M_{\mathbb A,\mathbb B,Y}(S_{\Phi,Y})
\]
is the total-response presentation of the same data whose uncurried form is
the primitive execution span of
\[
\mathsf{Prog}_{\Phi,Y}.
\]
Flattening the coalgebra recovers the span program category.  Pulling guarded
polynomial response liftings back along the coalgebra recovers the guarded span
modalities.  The terminal trace polynomial is the special case
\[
Y=\mathbf 1_{\mathbb B}.
\]
\end{remark}